% Paper B — draft of Sections 1-5, built from P5 (p5_v47), P4 (p4_v41), E2.
% House style: literal author-year prose citations, no citation macros.
% Inline maths uses \( ... \). External bibliography is NOT yet assembled;
% see the note at the end of this file.

\section{Introduction}
\label{sec:intro}

Environmental governance runs on clocks. Institutions observe a resource system at
discrete times, form a judgement from what they observe, and hold the resulting control
until the next review. The interval between reviews is not a detail of implementation: it
is a design variable, it varies across institutions, and it has stability consequences.
This paper is about that variable, and about the second clock that runs beside it --- the
lag from an observed decline to the institutional response that the observation was meant
to trigger. We call the pair the \emph{decision clock}.

The claim of this paper is that the decision clock, and not biology alone, decides whether
management stabilises or destabilises a renewable resource. That is a claim about a
mechanism, and it is worth stating at the outset what it is not. It is not a claim that
biological parameters are unimportant; the growth rate, the carrying capacity and the
shape of the surplus curve all enter what follows. It is a claim that, holding the biology
fixed, changing the timing and form of policy revision can move a system across a
stability boundary, and that this movement is large enough to matter at the review
intervals institutions actually use.

Two mechanisms make up the clock, and they are often conflated. The first is the
\emph{review interval} itself: a system managed by sample-and-hold --- observe, then hold
until the next review --- is a different dynamical object from the same system managed
continuously, and the substitution is not innocent. The second is the \emph{response
delay}: the time from observation to implemented response. The first is a property of the
institution's calendar; the second is a property of its deliberative and administrative
machinery. Both are design variables. Both have thresholds. They are analysed separately
in Sections \ref{sec:review} and \ref{sec:delay}, and the argument that they are two
halves of one idea rather than two independent effects is made in Section
\ref{sec:oneidea}.

Three lines of evidence support the claim, and they are reported in the order in which
they bear on it.

The first two are screens, and both are null. A multiplicity-controlled screen of
\textbf{42 fish stocks} finds no robust institutional cycles. A \textbf{32-system}
cross-sector case search finds no unconfounded oscillator. These null results are reported
here as results and not as limitations, for a reason that is substantive rather than
rhetorical: if periodicity alone could diagnose governance feedback, these screens would
have found the cycles, and they did not. The absence is what licenses the mechanism
argument. It says that the decision clock has to be identified structurally --- from the
form of the policy revision rule and the cadence on which it is applied --- rather than
inferred from the spectral signature of a management time series. Sections
\ref{sec:42} and \ref{sec:32} report them.

The third is a certification, and it is a single stock examined in depth. Northern cod
(NAFO divisions 2J3KL) is the case on which the certification apparatus is run, under a
protocol frozen before any score was computed. Section \ref{sec:cod} reports it. The
reason one stock can carry this weight is that the certificate is a characterisation
rather than a simulation: on the increasing branch of the surplus curve, the robustness of
a single-threshold reference point to persistent productivity shocks reduces to three
constants, and the kernel computation collapses to algebra. What is certified is therefore
a statement about the structure of the problem and not a statement about the fit.

A word on what this paper does \emph{not} claim, since the temptation is considerable and
the materials invite it. The review interval at which one of our systems loses stability
and the horizon over which the cod certificate holds are numerically close. That proximity
is a coincidence, it is an artefact of comparing two different systems, and no inference
is drawn from it anywhere in this paper. Section \ref{sec:conflict} explains why in
detail, and gives the computation that settles it. The short version is that the two
quantities point in opposite directions: longer review intervals stabilise one loop and
destroy the other. A paper that conflated them would be contradicted by both of the
results it rests on.

\section{The decision clock: two mechanisms}
\label{sec:clock}

\subsection{Review interval and sampled-control stability}
\label{sec:review}

Modelling periodic review as \emph{sample-and-hold governance} --- a loop that observes
and updates only at review times --- replaces the familiar continuous-time or single-step
annual formulations with a hybrid operator. Two properties of that operator are exact.
First, the sampled state space is forward invariant: the state computed at review times
remains in the set the construction is defined on. Second, rapid-review consistency holds
over finite horizons only: as the review interval shrinks the sampled system approaches
the continuous one on any fixed finite horizon, but the convergence is not uniform in the
horizon, so no limiting argument transfers a stability conclusion from the short-interval
limit to the intervals institutions actually use.

The consequence is that stability boundaries \emph{move or vanish when the operator is
changed}. This is the central technical fact of the review-interval mechanism, and it is
worth being precise about what it says. It does not say that approximations are
conservative, or that they are accurate at small intervals, or that they fail gracefully.
It says that the set of review intervals on which the system is stable is a property of
the operator, and that substituting a different operator --- a continuous delay, a
one-step annual map --- can relocate that set or empty it.

At the illustrative baseline used to demonstrate the mechanism, the exact map crosses the
stability boundary once, near a \textbf{6.5-year} review interval. One-step
approximations of the same system report crossings that are artefacts of the
discretisation. Under the protective rule the channel is stable at every interval tested,
with no crossing at all. So the boundary's location depends on the rule as well as on the
operator, and under one of the two rules there is no boundary to locate.

Two caveats attach to the 6.5-year figure and both travel with it throughout this paper.
The first is that the figure is \emph{ill-conditioned}: a two-hundredth-part change in the
exploitation ratio moves the crossing by more than half. The number is a property of one
knife-edge calibration, not a robust timescale of the mechanism, and it should be read
with a sensitivity band rather than as a point. The second caveat is specific to the
comparison with the cod certification and is deferred to Section \ref{sec:conflict}.

\subsection{Response delay and the stabilising window}
\label{sec:delay}

The second mechanism is the delay that sits between an observed decline and the
institutional response to it. This is not the ecological delay of recruitment or
maturation, and it is not the informational delay of a stock assessment catching up with
reality; both of those are well studied. It is \emph{governance delay}, the lag inside the
managing institution.

The model is a three-state system: a harvested stock, a filtered deficit memory, and
bounded effort. Two rules are compared. The \textbf{mobilising rule} has effort growing as
more effort is deployed --- a positive feedback in the response itself. The
\textbf{protective rule} is a quota-tracking law that restores harvest toward a cap.

The two rules behave in opposite ways, and the contrast is the result.

Under the \textbf{mobilising rule}, intermediate delay is \emph{stabilising}. There are
two subcritical Hopf crossings, interval-certified near \textbf{3.7} and \textbf{150}
years, and they bound a window of delay within which the lag itself stabilises the loop.
This is the counter-intuitive case and the one that matters for design: there is a range
of response lags that are not merely tolerable but load-bearing, and a manager who
shortens the lag out of that range destabilises a system that was stable. Delay is not a
defect to be minimised. It is a coordinate to be placed.

Under the \textbf{protective rule} the loop gain stays below one at the calibrated point,
and the equilibrium is exponentially stable at every delay. This is a \emph{no-Hopf
theorem}: there is no delay at which the loop destabilises, and consequently no window to
place. The apparent destabilisation that a discretised treatment of this rule reports near
a 2.3-year threshold is an artefact of that discretisation, the same failure mode that
produces the artefact crossings of Section \ref{sec:review}.

Under periodic review the two rules diverge again, and in the direction the continuous
analysis predicts. Annual protective review remains stable. Annual mobilising review is
unstable, and restabilises only above a 6.5-year interval, through a
Neimark--Sacker-type crossing in which a complex pair leaves the unit circle.

The attractor topology of the delay system is five-regime. Two folds are certified at the
discrete collocation level, with the continuum stages of that certification left open, and
there is a basin-boundary transition toward a large-amplitude attractor that is identified
but not verified. Those two open items are stated as open.

Finally, the model class is \emph{stock-agnostic}. Its design coordinates --- the sign of
the rule, the deployment delay, the review cadence --- are not properties of any
particular fishery, and they carry to any periodically reviewed renewable resource. That
is what licenses the cross-sector search of Section \ref{sec:32}.

\subsection{One idea, not two}
\label{sec:oneidea}

The review interval and the response delay are often treated as independent
imperfections --- one a matter of calendar, the other of administrative capacity. They are
not independent, and the reason is that both enter the same closed loop at the same place.

Both determine \emph{how stale the control is when it is applied}. Under sample-and-hold,
the control held between reviews is a function of a state that has moved on; the review
interval sets how far it has moved. Under governance delay, the response implements a
judgement about a state that has also moved on; the delay sets how far. In both cases the
loop contains a block whose output is a lagged function of its input, and in both cases
the stability question is whether that block's phase shift at the loop's crossover
frequency is small enough to preserve the loop's margin.

The mechanisms differ in where the lag is incurred and therefore in what can be changed
about it. The review interval is set by a calendar and can in principle be retimed
directly. The response delay is set by an institution's deliberative machinery and is
usually changed only indirectly, by changing what triggers a response or what
pre-authorised action follows it. But the two are coordinates of one design problem:
choose the cadence and the response law together, because the stability of the loop
depends on both and neither can be chosen well alone.

This is why the mechanisms are analysed separately and then recombined. Sections
\ref{sec:review} and \ref{sec:delay} establish what each does on its own. Section
\ref{sec:conflict} asks what happens when both are placed on one system, and that is where
the result of this paper lives.

\section{Evidence I: forty-two stocks, and the absence of institutional cycles}
\label{sec:42}

If the decision clock governs stability, the most direct evidence would be institutional
cycles in managed stocks: periodicity in a management time series at the review cadence,
or an oscillation whose period tracks the institution's calendar. A
multiplicity-controlled screen of \textbf{42 fish stocks} looks for exactly that, and
finds no robust institutional cycles.

The result is reported here as a result rather than as a limitation of the screen, and the
distinction matters. A screen that finds nothing can mean that nothing is there, or that
the instrument was too weak to find it. The screen was designed to separate those
readings: multiplicity control is applied precisely so that a null outcome can be
interpreted, rather than dismissed as an artefact of multiple comparisons. The finding is
that the periodicity is not there to be found.

What follows from the absence is the substantive point. Periodicity alone cannot diagnose
governance feedback. The spectral signature of a management time series does not reveal
the cadence that produced it, and the absence of an institutional cycle in a stock record
is not evidence that the decision clock is well set --- only that the clock leaves no
oscillatory fingerprint at the resolution the record supports. A well-governed stock and a
badly-governed stock can both be spectrally quiet.

This is what makes the decision clock a design variable that has to be reasoned about
structurally. It cannot be read off the data. It has to be placed, using the mechanism of
Section \ref{sec:clock} and the certification of Section \ref{sec:cod}, and then monitored
against the horizon it has to sit inside.

\section{Evidence II: thirty-two systems across sectors}
\label{sec:32}

The second screen asks whether the mechanism generalises beyond fisheries. A
\textbf{32-system} cross-sector case search looks for an unconfounded oscillator --- an
oscillation attributable to the governance loop rather than to the underlying system ---
and finds none.

The design logic is the same as in Section \ref{sec:42}, and the null has the same
positive content: across sectors whose review cadences, response laws and biological
timescales all differ, no case presents an oscillation that can be attributed to the
governance loop without confounding. The mechanism of Section \ref{sec:clock} is
stock-agnostic in its design coordinates, and the cross-sector search is the check on
whether that agnosticism is real or merely asserted. It is real: the coordinates carry,
and the absence of a diagnostic signature carries with them.

Taken together, the two screens say the same thing at two scales. Within fisheries and
across sectors, the decision clock does not announce itself in the data. Any claim about
it has to be made structurally, from the rule and the cadence, and certified against the
consequences rather than fitted to the record.

\section{The cod certification}
\label{sec:cod}

The third line of evidence is a single stock examined in full, under a protocol frozen
before any score was computed. The stock is Northern cod, NAFO divisions 2J3KL, with a
Schaefer fit over 1983--2007 and a limit reference point of \(884.6\) kt.

What makes one stock able to carry this weight is that the result is a
\emph{characterisation} rather than a simulation. On the increasing branch of the surplus
curve, the robustness of a single-threshold reference point to persistent productivity
shocks reduces to three constants, and the kernel computation collapses to algebra. A
certificate obtained this way is not a statement about the fit; it is a statement about
the structure of the problem, and it is exactly as strong as its assumptions.

\subsection{The three constants}

The first constant is the largest catch that holds the reference point from itself,
\[
C^{*} = g(\mathrm{LRP}) - |e| = 91.59\ \mathrm{kt}.
\]
Its significance is that harvest and protection are, at that point, one budget. Any rule's
protection margin is exactly \(C^{*}\) minus the catch it takes there. There is no
trade-off to be optimised and no curve to be traced: the margin is the constant minus the
catch.

The second constant is
\[
g_{\max} - |e| = 215.2\ \mathrm{kt},
\]
which separates viable-but-raised kernels from empty ones. Together the two constants
reproduce the kernel table in closed form, which is the check that the characterisation is
doing real work rather than restating the computation.

The third constant is the expansion rate of the map at the reference point. At the
calibrated point it is
\[
a_{\max} = F'(K^{*}) = 1.1531,
\]
and the map is expansive at the reference point for every admissible \(K \ge 2K^{*}\).

\subsection{Two consequences}

The first consequence is that the \textbf{no-dominance verdict is structural}. Because any
rule's protection margin is \(C^{*}\) minus its catch, a reactive rule cannot out-supply a
flat cap that it matches in protection. This is not an empirical finding about the rules
examined; it is an identity, and it holds for every rule in the class.

The second consequence is the one that matters for the decision clock, and it is the
reason this section belongs in this paper. Because the map is expansive at the reference
point, the certified layer has a \textbf{finite horizon}: \textbf{six years} under the two
harsher floors, \textbf{seven} under the informative one. That horizon is the crossing at
which a geometrically growing erosion margin overtakes the worst-case trajectory.

The decisive property of that horizon is that \textbf{no catch reduction extends it}. It is
the same for every admissible catch, including none at all. This is what makes it a
statement about the clock rather than about the harvest: the certificate does not expire
because too much is being taken, and no reduction in removals buys additional certified
time. It expires because the map is expansive at the reference point and the erosion
margin grows geometrically, and the only thing that changes the horizon is a change in
what is being certified or in the map itself.

One identification limit is stated plainly. Carrying capacity is not identified from
above; the functionals the results depend on are identified, and \(K\) is not. The results
are therefore conditional on the admissible set of \(K\), and every statement above is
made for every admissible \(K\), not for an estimated one.

% ---------------------------------------------------------------------------
% NOTE ON BIBLIOGRAPHY, to be resolved before this draft is committed.
%
% This draft deliberately contains almost no external citations rather than
% invented ones. All three source papers (p5_v47, p4_v41, paperE2_cod_intervention_v29)
% use LITERAL author-year prose citations and hand-written References sections,
% consistent with the house style already recorded for the obstruction
% manuscript. Paper B's bibliography should be assembled by merging the three
% source reference lists and de-duplicating, NOT by writing new entries.
%
% Every quantitative claim above is taken from the three source abstracts or
% from the settled computations recorded in FAMILY_CONSOLIDATION_PLAN.md
% sections 2 and 8.2. No number in this draft was estimated or rounded into
% place by the drafter.
% ---------------------------------------------------------------------------
